\section{CIC: The functional fragment}
\vspace{10.5cm}

We introduce a grammar of expressions of the Calculus of Construction.

It is indeed just a fragment of all possible expressions at the moment. We will expand the grammar in the next
sections part by part.

We also introduce a grammar for a context.

\begin{align*}
	uvar    {}::={} & u \vert uvar^\prime\\
	level   {}::={} & uvar \vert 0 \vert S uvar \vert \max(uvar, uvar) \vert imax(uvar, uvar)\\
	var     {}::={} & x \vert var^\prime\\
	expr    {}::={} &  var \vert U_{level} \vert expr expr \vert \lambda var:expr.expr \vert \forall var:expr.expr \\
	context {}::={} &  \diamond\vert context, var: expr 
\end{align*}

Terms of type $expr$ are called \textbf{expressions} and terms of type context \textbf{contexts}

Here are some examples of expressions:

\begin{align*}
	x\\
	U_{S0}\\
	\lambda x:U_{S0}.x\\
	\lambda x:U_{0}.x\lambda x^\prime:U_{S0}.x x^\prime\\
	\lambda x:U_{0}.x x^\prime\\
\end{align*}

Terms of kind $uvar$ are called universe variables. Informally, we write $u_0, u_1, u_2, v_0, v_1, \dots$ for these
variables.

The terms of kind $level$ are called \textbf{level labels}.

In place of $0, S0, SS0,\dots$ we will simply write $0, 1, 2,\dots$. Hence, all of the following are level labels:

$$u, v, u_0, 7, max(u, 5), imax(v, 7)$$

Analogously, all other lowercase latin letters apart from $u$ and $v$ are used as \textbf{expression variables}, i. e.
terms of type $var$.

